\documentclass[12pt]{article}
\usepackage[a4paper,margin=1in]{geometry}
\usepackage{amsmath, amssymb}
\usepackage{enumitem}
\usepackage{titlesec}

\title{\textbf{Lecture Scribe: Randomized Min-Cut Algorithm}}
\author{CSE400 --- Fundamentals of Probability in Computing}
\date{}

\begin{document}
\maketitle

\section{Min-Cut Problem}

\subsection{Why Use Min-Cut}
The min-cut algorithm is used in applications related to:
\begin{itemize}
    \item Network connectivity, reliability, and optimization
    \item Network design to improve communication efficiency and determine minimum capacity cuts
    \item Communication network robustness and fault tolerance
    \item VLSI circuit partitioning to reduce interconnect complexity
\end{itemize}

\subsection{Definition of Cut-Set}
A \textbf{cut-set} in a graph is a set of edges whose removal separates the graph into two or more connected components.

Given a graph $G=(V,E)$ with $n$ vertices, the \textbf{minimum cut (min-cut)} problem is to find a cut-set of minimum cardinality.

\subsection{Edge Contraction Operation}
The primary operation used in min-cut algorithms is \textbf{edge contraction}.

Contracting an edge $(u,v)$:
\begin{itemize}
    \item Merges vertices $u$ and $v$ into a single vertex
    \item Removes all edges between $u$ and $v$
    \item Retains all other edges
    \item May create parallel edges but no self-loops
\end{itemize}

\subsection{Successful and Unsuccessful Runs}
\textbf{Successful min-cut run:}  
An execution of a min-cut algorithm that correctly identifies the minimum cut.

\textbf{Unsuccessful min-cut run:}  
An execution where the algorithm fails to identify the true minimum cut.

\section{Max-Flow Min-Cut Theorem}

\subsection{Theorem Statement}
In a flow network, the \textbf{maximum flow} from source to sink equals the \textbf{total weight of edges in a minimum cut}.

\subsection{Associated Definitions}
\begin{itemize}
    \item \textbf{Capacity of a cut:} Sum of capacities of edges directed from set $X$ to set $Y$
    \item \textbf{Minimum cut:} Cut with smallest capacity
    \item \textbf{Maximum flow:} Largest feasible flow from source $S$ to sink $T$
\end{itemize}

\section{Deterministic Min-Cut Algorithm}

\subsection{Stoer--Wagner Theorem}
Let $s$ and $t$ be vertices of $G$, and let $G/\{s,t\}$ denote the graph formed by merging them.

A minimum cut of $G$ is the smaller of:
\begin{itemize}
    \item A minimum $s$--$t$ cut in $G$
    \item A minimum cut in $G/\{s,t\}$
\end{itemize}

This holds because:
\begin{itemize}
    \item Either a minimum cut separates $s$ and $t$
    \item Or no minimum cut separates them, so contraction preserves the minimum cut
\end{itemize}

\subsection{Pseudocode}

\textbf{Algorithm: MinimumCutPhase$(G,a)$}
\begin{enumerate}
    \item Initialize $A \leftarrow \{a\}$
    \item While $A \neq V$:
    \begin{itemize}
        \item Add the most tightly connected vertex to $A$
    \end{itemize}
    \item Return the cut weight of this phase
\end{enumerate}

\textbf{Algorithm: MinimumCut$(G)$}
\begin{enumerate}
    \item While $|V| \ge 1$:
    \begin{itemize}
        \item Choose any vertex $a$
        \item Execute MinimumCutPhase$(G,a)$
        \item If phase cut is lighter than current minimum, store it
        \item Shrink $G$ by merging the last two added vertices
    \end{itemize}
    \item Return the minimum cut
\end{enumerate}

\section{Randomized Min-Cut Algorithm}

\subsection{Motivation for Randomization}
Randomized algorithms:
\begin{itemize}
    \item Provide probabilistic success guarantees
    \item Improve efficiency and allow parallelization
    \item Offer approximation guarantees
    \item Avoid worst-case deterministic behavior
    \item Exhibit robustness and heuristic effectiveness
\end{itemize}

\subsection{Karger’s Randomized Algorithm}
Karger’s algorithm repeatedly:
\begin{itemize}
    \item Selects a random edge
    \item Contracts the edge
    \item Continues until only two supernodes remain
\end{itemize}

The remaining edges between the two supernodes define a cut, which may or may not be minimum.

\subsection{Recursive Randomized Min-Cut Pseudocode}

\textbf{Recursive-Randomized-Min-Cut$(G,\alpha)$}
\begin{enumerate}
    \item Input: Undirected multigraph $G$ with $n$ vertices and constant $\alpha>0$
    \item Output: A cut $C$ of $G$
    \item If $n \le \alpha^3$:
    \begin{itemize}
        \item Compute min-cut by exhaustive search
    \end{itemize}
    \item Else:
    \begin{itemize}
        \item Repeat for $i=1$ to $\alpha$:
        \begin{itemize}
            \item Obtain $G'$ by performing $n-\lfloor n/\sqrt{\alpha}\rfloor$ random contractions
            \item Compute $C' =$ Recursive-Randomized-Min-Cut$(G',\alpha)$
            \item If $i=1$ or $|C'|<|C|$, update $C \leftarrow C'$
        \end{itemize}
    \end{itemize}
    \item Return $C$
\end{enumerate}

\section{Deterministic vs Randomized Min-Cut}

\subsection{Deterministic Min-Cut}
\begin{itemize}
    \item Guarantees exact minimum cut
    \item Higher complexity for large graphs
    \item Stoer--Wagner complexity: $O(VE + V^2 \log V)$
\end{itemize}

\subsection{Randomized Min-Cut}
\begin{itemize}
    \item Produces minimum cut with high probability
    \item Time complexity: $O(V^2)$
\end{itemize}

\section{Probability of Correct Min-Cut}
The randomized algorithm outputs a true minimum cut with probability at least:
\[
\frac{2}{n(n-1)}.
\]

\section{Python Simulation}
Students are instructed to download the provided Jupyter notebook from the course platform and execute the simulation corresponding to Lecture~10.

\end{document}
